%% ============================================================
%%  SECCIÓN 2 — Espacios de Hilbert, Notación de Dirac y Operadores
%% ============================================================

\section{Espacios de Hilbert, Notación de Dirac y Operadores Físicos}
\label{sec:espacio_hilbert}

\subsection{El Espacio de Hilbert como Escenario Cuántico}

En la formulación rigurosa de von Neumann, el estado físico instantáneo de un sistema cuántico aislado queda descrito por un rayo vectorial (un vector normalizado salvo una fase global arbitraria $\eu^{\iu\alpha}$) que habita en un \textbf{espacio de Hilbert complejo}, denotado por $\mathcal{H}$ \cite{dirac, sakurai}.

\begin{teoremabox}[Definición Formal: Espacio de Hilbert $\mathcal{H}$]
    Un espacio de Hilbert $\mathcal{H}$ es un espacio vectorial lineal sobre el cuerpo de los números complejos $\mathbb{C}$, dotado de un producto interno $\braket{\cdot}{\cdot}: \mathcal{H} \times \mathcal{H} \to \mathbb{C}$, que induce una norma $\norm{\psi} = \sqrt{\braket{\psi}{\psi}}$, y que es \textbf{completo} respecto a la métrica inducida por dicha norma (es decir, toda sucesión de Cauchy converge a un elemento dentro del propio espacio).
\end{teoremabox}

En mecánica cuántica no relativista de una partícula en tres dimensiones, $\mathcal{H}$ suele identificarse con el espacio funcional $\mathcal{L}^2(\mathbb{R}^3)$ de funciones de cuadrado integrable con respecto a la medida de Lebesgue:
\begin{equation}
    \int_{\mathbb{R}^3} |\psi(\vec{r})|^2\,\dif^3\vec{r} < \infty.
    \label{eq:l2_norm}
\end{equation}
Para grados de libertad discretos (como el espín $s = 1/2$), $\mathcal{H}$ es un espacio de dimensión finita isomorfo a $\mathbb{C}^2$.

\subsection{Dualidad Bra-Ket y Axiomática del Producto Interno}

Paul Dirac introdujo una notación sumamente elegante y potente que desacopla la física de cualquier sistema particular de coordenadas:
\begin{itemize}
    \item Los elementos de $\mathcal{H}$ se denominan \textbf{kets} y se denotan por $\ket{\psi}$.
    \item Los elementos del espacio dual continuo $\mathcal{H}^*$ (el espacio de funcionales lineales continuos que transforman kets en números complejos $\mathbb{C}$) se denominan \textbf{bras} y se denotan por $\bra{\phi}$.
\end{itemize}

Por el \textit{Teorema de Representación de Riesz}, existe un isomorfismo antilineal canónico entre $\mathcal{H}$ y su dual $\mathcal{H}^*$: a todo ket $\ket{\psi}$ le corresponde biunívocamente un bra $\bra{\psi}$, denominado su adjunto o conjugado hermítico:
\begin{equation}
    \ket{\psi}\adj = \bra{\psi}, \qquad (\alpha\ket{\phi} + \beta\ket{\psi})\adj = \alpha^*\bra{\phi} + \beta^*\bra{\psi}.
    \label{eq:riesz_adjoint}
\end{equation}

El producto interno entre un bra $\bra{\phi}$ y un ket $\ket{\psi}$ se denota por el \textit{bracket} $\braket{\phi}{\psi} \in \mathbb{C}$ y satisface los siguientes axiomas fundamentales:
\begin{enumerate}
    \item \textbf{Linealidad en el segundo argumento (ket):}
    \begin{equation}
        \braket{\phi}{\alpha\psi_1 + \beta\psi_2} = \alpha\braket{\phi}{\psi_1} + \beta\braket{\phi}{\psi_2}, \quad \forall\,\alpha, \beta \in \mathbb{C}.
    \end{equation}
    \item \textbf{Antilinealidad en el primer argumento (bra):}
    \begin{equation}
        \braket{\alpha\phi_1 + \beta\phi_2}{\psi} = \alpha^*\braket{\phi_1}{\psi} + \beta^*\braket{\phi_2}{\psi}.
    \end{equation}
    \item \textbf{Simetría conjugada hermítica:}
    \begin{equation}
        \braket{\phi}{\psi} = \braket{\psi}{\phi}^*.
        \label{eq:hermitian_symmetry}
    \end{equation}
    \item \textbf{Positividad definida:}
    \begin{equation}
        \braket{\psi}{\psi} \ge 0, \quad \text{y} \quad \braket{\psi}{\psi} = 0 \iff \ket{\psi} = 0.
        \label{eq:positive_definite}
    \end{equation}
\end{enumerate}

La condición física de normalización para un estado cuántico que representa una certeza de probabilidad total es:
\begin{equation}
    \norm{\psi}^2 = \braket{\psi}{\psi} = 1.
\end{equation}

\subsection{Bases Ortonormales y la Relación de Clausura}

Si $\{\ket{n}\}$ constituye una base discreta y ortonormal de $\mathcal{H}$, sus elementos satisfacen la condición de ortonormalidad:
\begin{equation}
    \braket{m}{n} = \delta_{mn},
    \label{eq:orthonormality}
\end{equation}
donde $\delta_{mn}$ es el delta de Kronecker. Cualquier estado arbitrario $\ket{\psi}$ puede expandirse de manera única como una combinación lineal:
\begin{equation}
    \ket{\psi} = \sum_n c_n \ket{n}.
\end{equation}
Multiplicando por la izquierda por el bra $\bra{m}$ y aplicando la Ec.~\eqref{eq:orthonormality}:
\begin{equation}
    \braket{m}{\psi} = \sum_n c_n \braket{m}{n} = \sum_n c_n \delta_{mn} = c_m.
\end{equation}
Sustituyendo estos coeficientes $c_n = \braket{n}{\psi}$ nuevamente en la expansión de $\ket{\psi}$:
\begin{equation}
    \ket{\psi} = \sum_n \ket{n}\braket{n}{\psi} = \left( \sum_n \ket{n}\bra{n} \right) \ket{\psi}.
\end{equation}
Como esta igualdad se cumple para cualquier vector $\ket{\psi} \in \mathcal{H}$, el operador entre paréntesis debe ser idéntico al operador identidad $\hat{I}$. Esta es la célebre \textbf{Relación de Clausura (o completitud)}:
\begin{equation}
    \boxed{\sum_n \ket{n}\bra{n} = \hat{I}}.
    \label{eq:closure_relation}
\end{equation}
La relación de clausura es una de las herramientas de cálculo más poderosas de la física cuántica: permite \enquote{interpolar} una base en cualquier producto interno o elemento de matriz, transformando ecuaciones abstractas de operadores en sumas algebraicas directas \cite{zettili}.

\subsection{Operadores Lineales y el Operador Adjunto Hermítico}

Un operador lineal $\hat{A}$ transforma vectores de $\mathcal{H}$ en otros vectores del mismo espacio:
\begin{equation}
    \hat{A}(\alpha\ket{\psi_1} + \beta\ket{\psi_2}) = \alpha\hat{A}\ket{\psi_1} + \beta\hat{A}\ket{\psi_2}.
\end{equation}
El \textbf{adjunto hermítico} de $\hat{A}$, denotado por $\hat{A}\adj$, se define formalmente a través del producto interno como el único operador que satisface:
\begin{equation}
    \braket{\phi}{\hat{A}\psi} = \braket{\hat{A}\adj\phi}{\psi}, \quad \forall\,\ket{\phi}, \ket{\psi} \in \mathcal{H}.
    \label{eq:def_adjoint}
\end{equation}
Equivalentemente, en la notación de brackets, si $\ket{\chi} = \hat{A}\ket{\psi}$, entonces el bra asociado es $\bra{\chi} = \bra{\psi}\hat{A}\adj$. De esta definición se deducen directamente las siguientes propiedades esenciales:
\begin{enumerate}
    \item $(\hat{A} + \hat{B})\adj = \hat{A}\adj + \hat{B}\adj$.
    \item $(\alpha\hat{A})\adj = \alpha^* \hat{A}\adj$, con $\alpha \in \mathbb{C}$.
    \item Inversión de orden en productos de operadores:
    \begin{equation}
        (\hat{A}\hat{B})\adj = \hat{B}\adj \hat{A}\adj.
        \label{eq:adjoint_product}
    \end{equation}
    \textit{Demostración:} Aplicando la definición dos veces sucesivas,
    $\braket{\phi}{(\hat{A}\hat{B})\psi} = \braket{\hat{A}\adj\phi}{\hat{B}\psi} = \braket{\hat{B}\adj(\hat{A}\adj\phi)}{\psi} = \braket{(\hat{B}\adj\hat{A}\adj)\phi}{\psi}$.
    \item Involución: $(\hat{A}\adj)\adj = \hat{A}$.
\end{enumerate}

\begin{teoremabox}[Clasificación de Operadores según su Adjunción]
    \begin{itemize}
        \item \textbf{Operador Hermítico (Autoadjunto):} Cumple $\hat{A}\adj = \hat{A}$.
        En mecánica cuántica, \textbf{todo observable físico medible se representa por un operador lineal hermítico}. El Teorema Espectral garantiza que todos sus autovalores son estrictamente reales ($\lambda_n \in \mathbb{R}$) y que sus autovectores correspondientes a autovalores distintos son mutuamente ortogonales.
        \item \textbf{Operador Antihermítico (o Antiautoadjunto):} Cumple $\hat{C}\adj = -\hat{C}$.
        Sus autovalores son imaginarios puros o cero, y su valor esperado en cualquier estado físico es estrictamente imaginario puro: $\ev{\hat{C}} \in \iu\mathbb{R}$.
        \item \textbf{Operador Unitario:} Cumple $\hat{U}\adj\hat{U} = \hat{U}\hat{U}\adj = \hat{I}$, lo que equivale a $\hat{U}\adj = \hat{U}^{-1}$. Los operadores unitarios preservan productos internos y normas, describiendo simetrías y la evolución temporal continua del sistema.
    \end{itemize}
\end{teoremabox}

\subsubsection{Diccionario Matricial del Bra-Ket y el Producto Externo}
\label{subsec:diccionario_matricial_braket}

Para conectar sin fisuras la notación abstracta de Dirac con el cálculo matricial concreto de Álgebra Lineal, fijemos una base ortonormal discreta $\{\ket{1}, \ket{2}, \dots, \ket{N}\}$ de un espacio de Hilbert $\mathcal{H}$ de dimensión $N$.

\begin{enumerate}
    \item \textbf{Kets como Vectores Columna y Bras como Vectores Fila Conjugados:}
    Todo ket $\ket{\psi} = \sum_{n=1}^N c_n\ket{n}$ queda representado unívocamente por el vector columna de sus componentes $c_n = \braket{n}{\psi}$, mientras que su bra adjunto $\bra{\psi} = \sum_{n=1}^N c_n^*\bra{n}$ corresponde al vector fila transpuesto conjugado:
    \begin{equation}
        \ket{\psi} \;\doteq\; \begin{pmatrix} c_1 \\ c_2 \\ \vdots \\ c_N \end{pmatrix}, \qquad
        \bra{\psi} = \ket{\psi}\adj \;\doteq\; \begin{pmatrix} c_1^* & c_2^* & \cdots & c_N^* \end{pmatrix}.
    \end{equation}

    \item \textbf{Producto Interno $\braket{\phi}{\psi}$ frente a Producto Externo $\ket{\psi}\bra{\phi}$:}
    En la notación de Dirac, el orden en que se yuxtaponen un bra y un ket cambia radicalmente la naturaleza matemática del objeto:
    \begin{itemize}
        \item \textbf{Producto Interno (Bra-Ket $\braket{\phi}{\psi}$):} Es una fila $1\times N$ multiplicada por una columna $N\times 1$, lo que arroja un \textbf{número complejo escalar} ($\mathbb{C}$):
        \begin{equation}
            \braket{\phi}{\psi} = \begin{pmatrix} d_1^* & d_2^* & \cdots & d_N^* \end{pmatrix} \begin{pmatrix} c_1 \\ c_2 \\ \vdots \\ c_N \end{pmatrix} = \sum_{n=1}^N d_n^* c_n \in \mathbb{C}.
        \end{equation}
        \item \textbf{Producto Externo (Ket-Bra $\ket{\psi}\bra{\phi}$):} Es una columna $N\times 1$ multiplicada por una fila $1\times N$, lo que produce un \textbf{operador lineal} representado por una \textbf{matriz cuadrada de $N\times N$}:
        \begin{equation}
            \ket{\psi}\bra{\phi} \;\doteq\; \begin{pmatrix} c_1 \\ c_2 \\ \vdots \\ c_N \end{pmatrix} \begin{pmatrix} d_1^* & d_2^* & \cdots & d_N^* \end{pmatrix} = \begin{pmatrix}
                c_1 d_1^* & c_1 d_2^* & \cdots & c_1 d_N^* \\
                c_2 d_1^* & c_2 d_2^* & \cdots & c_2 d_N^* \\
                \vdots & \vdots & \ddots & \vdots \\
                c_N d_1^* & c_N d_2^* & \cdots & c_N d_N^*
            \end{pmatrix}.
        \end{equation}
        En efecto, al actuar sobre cualquier ket $\ket{\chi}$, el operador $(\ket{\psi}\bra{\phi})$ produce $\left(\ket{\psi}\bra{\phi}\right)\ket{\chi} = \ket{\psi}\braket{\phi}{\chi} = (\braket{\phi}{\chi})\ket{\psi}$, que es un vector en la dirección de $\ket{\psi}$.
    \end{itemize}

    \item \textbf{¿Cómo se construye la matriz de un operador $\hat{A}$ en una base dada?}
    Insertando dos veces la relación de clausura $\hat{I} = \sum_m \ket{m}\bra{m}$ a ambos lados del operador $\hat{A}$:
    \begin{equation}
        \hat{A} = \hat{I}\hat{A}\hat{I} = \left(\sum_{m=1}^N \ket{m}\bra{m}\right) \hat{A} \left(\sum_{n=1}^N \ket{n}\bra{n}\right) = \sum_{m=1}^N \sum_{n=1}^N \underbrace{\matrixel{m}{\hat{A}}{n}}_{A_{mn}} \ket{m}\bra{n}.
    \end{equation}
    El número complejo $A_{mn} \equiv \matrixel{m}{\hat{A}}{n}$ es exactamente el \textbf{elemento de matriz en la fila $m$ y columna $n$}:
    \begin{equation}
        \boxed{A_{mn} = \matrixel{m}{\hat{A}}{n} = \braket{m}{\hat{A}n} \quad (\text{fila } m, \text{ columna } n)}.
    \end{equation}
    \textit{Regla práctica para construir la matriz:} Aplique el operador $\hat{A}$ sobre el $n$-ésimo vector base $\ket{n}$; los coeficientes del vector resultante constituyen la \textbf{$n$-ésima columna} de la matriz $A$.
\end{enumerate}

\begin{table}[H]
    \centering
    \caption{Diccionario de traducción entre el Álgebra Lineal Matricial y la Notación de Dirac.}
    \label{tab:diccionario_dirac}
    \vspace{0.15cm}
    \resizebox{\linewidth}{!}{%
    \begin{tabular}{@{}lll@{}}
        \toprule
        \textbf{Concepto Matemático} & \textbf{Notación de Dirac (Abstracta)} & \textbf{Representación Matricial en Base $\{\ket{n}\}$} \\
        \midrule
        Vector de estado & Ket $\ket{\psi}$ & Vector columna $N\times 1$: $(c_1, c_2, \dots, c_N)^T$ \\
        Funcional dual (adjunto) & Bra $\bra{\psi} = \ket{\psi}\adj$ & Vector fila conjugado $1\times N$: $(c_1^*, c_2^*, \dots, c_N^*)$ \\
        Producto interno (escalar) & Bracket $\braket{\phi}{\psi}$ & Producto fila por columna: $\sum_n d_n^* c_n \in \mathbb{C}$ \\
        Producto externo (operador) & Ket-Bra $\ket{\psi}\bra{\phi}$ & Matriz cuadrada $N\times N$ de entradas $c_m d_n^*$ \\
        Elemento de matriz de $\hat{A}$ & Sandwich $\matrixel{m}{\hat{A}}{n}$ & Entrada $A_{mn}$ (fila $m$, columna $n$) \\
        Operador adjunto $\hat{A}\adj$ & $\matrixel{m}{\hat{A}\adj}{n} = \matrixel{n}{\hat{A}}{m}^*$ & Matriz transpuesta conjugada: $(A\adj)_{mn} = A_{nm}^*$ \\
        Traza de un operador & $\mathrm{Tr}(\hat{A}) = \sum_n \matrixel{n}{\hat{A}}{n}$ & Suma de la diagonal principal: $\sum_n A_{nn}$ \\
        \bottomrule
    \end{tabular}%
    }
\end{table}

\subsection{Valor Esperado, Varianza y el Operador Desviación}

Cuando un sistema se encuentra en un estado físico normalizado $\ket{\psi}$ ($\braket{\psi}{\psi}=1$), la medición de un observable $\hat{A}$ no arroja en general un valor determinista, sino una distribución estadística de resultados.

\subsubsection{El Teorema Espectral y el Postulado de Medida de von Neumann}
¿De dónde surge la fórmula $\ev{\hat{A}} = \matrixel{\psi}{\hat{A}}{\psi}$? No es una definición ad-hoc, sino una consecuencia directa del \textbf{Teorema Espectral para operadores autoadjuntos} en espacios de Hilbert.

Dado un observable hermítico $\hat{A}$, existe una base ortonormal completa de autovectores $\{\ket{n}\}$ con autovalores reales $\{a_n\}$ tales que $\hat{A}\ket{n} = a_n \ket{n}$. Utilizando la relación de clausura $\sum_n \ket{n}\bra{n} = \hat{I}$, el operador $\hat{A}$ admite la descomposición espectral única en proyectores ortogonales unidimensionales $\hat{P}_n \equiv \ket{n}\bra{n}$:
\begin{equation}
    \hat{A} = \hat{A}\hat{I} = \hat{A}\left( \sum_n \ket{n}\bra{n} \right) = \sum_n a_n \ket{n}\bra{n} = \sum_n a_n \hat{P}_n.
    \label{eq:spectral_decomposition}
\end{equation}
donde los proyectores satisfacen $\hat{P}_n^2 = \hat{P}_n$, $\hat{P}_n\adj = \hat{P}_n$ y $\hat{P}_n\hat{P}_m = \delta_{nm}\hat{P}_n$.

Cuando el sistema está en un estado normalizado arbitrario $\ket{\psi}$, podemos expandirlo en dicha base ortonormal:
\begin{equation}
    \ket{\psi} = \sum_n c_n \ket{n}, \quad c_n = \braket{n}{\psi}.
\end{equation}
Por el \textbf{Postulado de Medida de Born y von Neumann}:
\begin{enumerate}
    \item Los únicos resultados posibles al medir el observable $\hat{A}$ son sus autovalores $\{a_n\}$.
    \item La probabilidad de obtener el resultado discreto $a_n$ está dada por el valor esperado del proyector $\hat{P}_n$:
    \begin{equation}
        P(a_n) = \matrixel{\psi}{\hat{P}_n}{\psi} = \braket{\psi}{n}\braket{n}{\psi} = |c_n|^2.
    \end{equation}
    Como $\norm{\psi}^2 = 1$, la probabilidad total está conservada: $\sum_n P(a_n) = \sum_n |c_n|^2 = 1$.
    \item \textbf{Colapso del Estado (Postulado de Proyección de von Neumann):} Si al efectuar la medición sobre $\ket{\psi}$ el aparato registra el autovalor específico $a_n$, el estado físico del sistema sufre un cambio abrupto e irreversible, proyectándose instantáneamente sobre el subespacio propio correspondiente:
    \begin{equation}
        \ket{\psi} \;\xrightarrow{\quad\text{medición de } a_n\quad}\; \frac{\hat{P}_n\ket{\psi}}{\sqrt{\matrixel{\psi}{\hat{P}_n}{\psi}}} = \frac{c_n}{|c_n|}\ket{n} \;\sim\; \ket{n}.
    \end{equation}
    Esto garantiza la reproducibilidad física: si se vuelve a medir el mismo observable $\hat{A}$ inmediatamente después sobre el sistema recién colapsado, se obtendrá con certeza absoluta ($100\%$) el mismo valor $a_n$.
\end{enumerate}

\subsubsection{Valor Esperado}
El valor esperado $\ev{\hat{A}}$ es el promedio estadístico ponderado de todos los resultados posibles en un ensamble estadístico de mediciones idénticas:
\begin{align}
    \ev{\hat{A}} &\equiv \sum_n a_n P(a_n) = \sum_n a_n \braket{\psi}{n}\braket{n}{\psi} \nonumber\\
    &= \bra{\psi} \left( \sum_n a_n \ket{n}\bra{n} \right) \ket{\psi} = \boxed{\matrixel{\psi}{\hat{A}}{\psi}}.
    \label{eq:expectation_value}
\end{align}
¡La fórmula $\matrixel{\psi}{\hat{A}}{\psi}$ es exactamente el promedio probabilístico ponderado por el espectro del operador!

Si $\hat{A}$ es hermítico, su valor esperado es obligatoriamente un número real:
\begin{equation}
    \ev{\hat{A}}^* = \left(\matrixel{\psi}{\hat{A}}{\psi}\right)^* = \matrixel{\psi}{\hat{A}\adj}{\psi} = \matrixel{\psi}{\hat{A}}{\psi} = \ev{\hat{A}} \implies \ev{\hat{A}} \in \mathbb{R}.
\end{equation}

\subsubsection{El Operador Desviación \texorpdfstring{$\Delta\hat{A}$}{Delta A}}
Para cuantificar la dispersión de las medidas alrededor de la media, definimos el \textbf{operador desviación} (o fluctuación cuántica):
\begin{equation}
    \boxed{\Delta\hat{A} \equiv \hat{A} - \ev{\hat{A}}\hat{I}}.
    \label{eq:deviation_operator}
\end{equation}
Observemos con cuidado la naturaleza matemática de este operador: $\hat{A}$ es un operador, mientras que $\ev{\hat{A}}$ es un escalar (un c-número real). Para poder restarlos con sentido algebraico riguroso, multiplicamos formalmente el escalar por el operador identidad $\hat{I}$ (frecuentemente omitido en la notación abreviada, pero implícito en todo cálculo).

\begin{alertabox}[Propiedad Fundamental del Operador Desviación]
    Como $\hat{A}$ es hermítico y $\ev{\hat{A}} \in \mathbb{R}$, el operador desviación es automáticamente \textbf{hermítico}:
    \begin{equation}
        (\Delta\hat{A})\adj = (\hat{A} - \ev{\hat{A}}\hat{I})\adj = \hat{A}\adj - \ev{\hat{A}}^*\hat{I}\adj = \hat{A} - \ev{\hat{A}}\hat{I} = \Delta\hat{A}.
    \end{equation}
    Además, su valor esperado promedio es idénticamente nulo:
    \begin{equation}
        \ev{\Delta\hat{A}} = \ev{\hat{A} - \ev{\hat{A}}\hat{I}} = \ev{\hat{A}} - \ev{\hat{A}}\ev{\hat{I}} = \ev{\hat{A}} - \ev{\hat{A}} = 0.
    \end{equation}
\end{alertabox}

\subsubsection{Varianza Cuántica}
La incertidumbre cuántica (o desviación estándar) $\Delta A$ se define como la raíz cuadrada de la varianza:
\begin{equation}
    (\Delta A)^2 \equiv \ev{(\Delta\hat{A})^2} = \matrixel{\psi}{(\Delta\hat{A})^2}{\psi}.
    \label{eq:variance_def}
\end{equation}
Expandiendo el cuadrado del operador:
\begin{align}
    (\Delta\hat{A})^2 &= (\hat{A} - \ev{\hat{A}}\hat{I})(\hat{A} - \ev{\hat{A}}\hat{I}) \nonumber\\
    &= \hat{A}^2 - \hat{A}\ev{\hat{A}}\hat{I} - \ev{\hat{A}}\hat{I}\hat{A} + \ev{\hat{A}}^2 \hat{I}^2 \nonumber\\
    &= \hat{A}^2 - 2\ev{\hat{A}}\hat{A} + \ev{\hat{A}}^2 \hat{I}.
\end{align}
Tomando el valor esperado en el estado $\ket{\psi}$:
\begin{align}
    (\Delta A)^2 &= \ev{\hat{A}^2 - 2\ev{\hat{A}}\hat{A} + \ev{\hat{A}}^2 \hat{I}} \nonumber\\
    &= \ev{\hat{A}^2} - 2\ev{\hat{A}}\ev{\hat{A}} + \ev{\hat{A}}^2 = \boxed{\ev{\hat{A}^2} - \ev{\hat{A}}^2}.
    \label{eq:variance_expanded}
\end{align}
Esta fórmula es formalmente análoga a la varianza estadística clásica $\sigma_A^2 = \overline{A^2} - (\overline{A})^2$, pero con una diferencia conceptual titánica: la varianza cuántica no proviene de imperfecciones o ignorancia instrumental, sino de la dispersión ontológica intrínseca del estado cuántico $\ket{\psi}$.

\subsection{Ejemplo Resuelto 1: Cálculo de Dispersión Cuántica en un Sistema de Dos Niveles}

\begin{ejemplobox}[Ejemplo Resuelto 1: Fluctuación Cuántica en una Matriz de Pauli]
    \textbf{Enunciado:} Considere un sistema de dos niveles (espín $1/2$) con base ortonormal $\{\ket{+z}, \ket{-z}\}$. El observable físico $\hat{S}_x$ está dado por la matriz de Pauli:
    \begin{equation}
        \hat{S}_x = \frac{\hbar}{2}\begin{pmatrix} 0 & 1 \\ 1 & 0 \end{pmatrix}.
    \end{equation}
    Suponga que el sistema se encuentra preparado en el estado cuántico $\ket{\psi} = \frac{1}{\sqrt{5}}\begin{pmatrix} 2 \\ 1 \end{pmatrix}$.
    \begin{enumerate}
        \item Verifique que $\ket{\psi}$ está normalizado y calcule $\ev{\hat{S}_x}$.
        \item Calcule $\ev{\hat{S}_x^2}$ y determine la incertidumbre $\Delta S_x$.
        \item Construya explícitamente el operador desviación $\Delta\hat{S}_x$ en forma matricial y verifique directamente que $\matrixel{\psi}{(\Delta\hat{S}_x)^2}{\psi} = (\Delta S_x)^2$.
    \end{enumerate}

    \textbf{Solución Paso a Paso:}
    
    \textit{Paso 1: Normalización y Valor Esperado:}
    Calculamos el producto interno:
    \begin{equation}
        \braket{\psi}{\psi} = \frac{1}{\sqrt{5}}\begin{pmatrix} 2 & 1 \end{pmatrix} \frac{1}{\sqrt{5}}\begin{pmatrix} 2 \\ 1 \end{pmatrix} = \frac{1}{5}(2^2 + 1^2) = \frac{5}{5} = 1.
    \end{equation}
    El estado está correctamente normalizado. Evaluamos ahora la acción de $\hat{S}_x$ sobre $\ket{\psi}$:
    \begin{equation}
        \hat{S}_x \ket{\psi} = \frac{\hbar}{2}\begin{pmatrix} 0 & 1 \\ 1 & 0 \end{pmatrix} \frac{1}{\sqrt{5}}\begin{pmatrix} 2 \\ 1 \end{pmatrix} = \frac{\hbar}{2\sqrt{5}}\begin{pmatrix} 1 \\ 2 \end{pmatrix}.
    \end{equation}
    Tomando el producto interno con $\bra{\psi}$:
    \begin{equation}
        \ev{\hat{S}_x} = \braket{\psi}{\hat{S}_x\psi} = \frac{1}{\sqrt{5}}\begin{pmatrix} 2 & 1 \end{pmatrix} \frac{\hbar}{2\sqrt{5}}\begin{pmatrix} 1 \\ 2 \end{pmatrix} = \frac{\hbar}{10}(2\cdot 1 + 1\cdot 2) = \frac{4\hbar}{10} = \frac{2}{5}\hbar = 0{,}40\,\hbar.
    \end{equation}

    \textit{Paso 2: Varianza Cuántica:}
    Notemos que para las matrices de Pauli $\sigma_x^2 = \hat{I}_2$. Por ende:
    \begin{equation}
        \hat{S}_x^2 = \left(\frac{\hbar}{2}\right)^2 \begin{pmatrix} 0 & 1 \\ 1 & 0 \end{pmatrix}\begin{pmatrix} 0 & 1 \\ 1 & 0 \end{pmatrix} = \frac{\hbar^2}{4}\begin{pmatrix} 1 & 0 \\ 0 & 1 \end{pmatrix} = \frac{\hbar^2}{4}\hat{I}.
    \end{equation}
    En consecuencia, $\ev{\hat{S}_x^2} = \frac{\hbar^2}{4}\ev{\hat{I}} = \frac{\hbar^2}{4}$. Aplicando la fórmula de la varianza:
    \begin{equation}
        (\Delta S_x)^2 = \ev{\hat{S}_x^2} - \ev{\hat{S}_x}^2 = \frac{\hbar^2}{4} - \left(\frac{2}{5}\hbar\right)^2 = \frac{\hbar^2}{4} - \frac{4}{25}\hbar^2 = \frac{25 - 16}{100}\hbar^2 = \frac{9}{100}\hbar^2.
    \end{equation}
    Extrayendo la raíz cuadrada:
    \begin{equation}
        \Delta S_x = \sqrt{\frac{9}{100}\hbar^2} = \frac{3}{10}\hbar = 0{,}30\,\hbar.
    \end{equation}

    \textit{Paso 3: Matriz del Operador Desviación:}
    Por definición, $\Delta\hat{S}_x = \hat{S}_x - \ev{\hat{S}_x}\hat{I}$:
    \begin{equation}
        \Delta\hat{S}_x = \frac{\hbar}{2}\begin{pmatrix} 0 & 1 \\ 1 & 0 \end{pmatrix} - \frac{2}{5}\hbar\begin{pmatrix} 1 & 0 \\ 0 & 1 \end{pmatrix} = \hbar \begin{pmatrix} -2/5 & 1/2 \\ 1/2 & -2/5 \end{pmatrix} = \frac{\hbar}{10}\begin{pmatrix} -4 & 5 \\ 5 & -4 \end{pmatrix}.
    \end{equation}
    Actuando con este operador sobre $\ket{\psi}$:
    \begin{equation}
        \Delta\hat{S}_x \ket{\psi} = \frac{\hbar}{10\sqrt{5}}\begin{pmatrix} -4 & 5 \\ 5 & -4 \end{pmatrix}\begin{pmatrix} 2 \\ 1 \end{pmatrix} = \frac{\hbar}{10\sqrt{5}}\begin{pmatrix} -8 + 5 \\ 10 - 4 \end{pmatrix} = \frac{\hbar}{10\sqrt{5}}\begin{pmatrix} -3 \\ 6 \end{pmatrix}.
    \end{equation}
    Calculando la norma cuadrada de este vector resultante:
    \begin{equation}
        \matrixel{\psi}{(\Delta\hat{S}_x)^2}{\psi} = \norm{\Delta\hat{S}_x\psi}^2 = \frac{\hbar^2}{100\cdot 5}\left( (-3)^2 + 6^2 \right) = \frac{\hbar^2}{500}(9 + 36) = \frac{45}{500}\hbar^2 = \frac{9}{100}\hbar^2,
    \end{equation}
    lo cual coincide exactamente con el valor calculado en el Paso 2.
\end{ejemplobox}

\begin{figure}[H]
\centering
\begin{subfigure}[t]{0.50\textwidth}
\centering
\begin{tikzpicture}[scale=2.0,
   proj/.style={x={(1cm,0cm)},y={(0.45cm,0.3cm)},z={(0cm,1cm)}}]
  \draw[gray!70] (0,0) circle (1);
  \begin{scope}[proj]
    \draw[gray!60,dashed,domain=0:360,samples=100,smooth,variable=\t]
      plot ({cos(\t)},{sin(\t)},0);
    \draw[gray!60,dashed,domain=0:360,samples=100,smooth,variable=\t]
      plot ({cos(\t)},0,{sin(\t)});
    \draw[->,gray] (0,0,0) -- (1.35,0,0) node[right] {$x$};
    \draw[->,gray] (0,0,0) -- (0,1.35,0) node[right] {$y$};
    \draw[->,gray] (0,0,0) -- (0,0,1.45) node[above] {$z$};
    \node[font=\scriptsize,below] at (0,0,-1.0) {$\ket{{-}z}$};
    \node[font=\scriptsize,above right] at (0,0,1.0) {$\ket{{+}z}$};
    \node[font=\scriptsize,unigreen,below right] at (1.0,0,0) {$\ket{{+}x}$};
    \draw[dotted,unired] (0.8,0,0.6) -- (0.8,0,0);
    \draw[dotted,unired] (0.8,0,0) -- (0,0,0);
    \draw[-{Stealth},very thick,unired] (0,0,0) -- (0.8,0,0.6);
    \fill[unired] (0.8,0,0.6) circle (0.025);
    \node[unired,font=\small,right] at (0.8,0,0.6) {$\vec r$};
  \end{scope}
\end{tikzpicture}
\caption{$\vec r=(0{,}8;\,0;\,0{,}6)$, con $\langle S_x\rangle=\tfrac{\hbar}{2}r_x=0{,}4\,\hbar$.}
\end{subfigure}\hfill
\begin{subfigure}[t]{0.46\textwidth}
\centering
\begin{tikzpicture}
\begin{axis}[ybar,width=\linewidth,height=5.6cm,ymin=0,ymax=1.1,
  bar width=22pt,symbolic x coords={m,p},xtick=data,
  xticklabels={$-\hbar/2$,$+\hbar/2$},
  ylabel={Probabilidad},xlabel={Resultado de $S_x$},
  nodes near coords,nodes near coords align={vertical},
  enlarge x limits=0.5]
  \addplot[fill=uniblue!60,draw=uniblue] coordinates {(m,0.1) (p,0.9)};
\end{axis}
\end{tikzpicture}
\caption{$P(\pm)=\frac{1\pm0{,}8}{2}$; $\Delta S_x=0{,}3\,\hbar$.}
\end{subfigure}
\caption{Ejemplo Resuelto 1: estado $\ket\psi=(2,1)^T/\sqrt5$ sobre la esfera de Bloch y distribuci\'on de resultados al medir $\hat S_x$.}
\label{fig:bloch}
\end{figure}

\subsection{Teorema Espectral y Completitud (Teorema 2.2)}
\label{subsec:teorema_2_2}

En el estudio del álgebra lineal cuántica, la descomposición espectral representa el puente definitivo entre los operadores lineales y la física de las mediciones experimentales. El postulado de que los observables físicos son operadores autoadjuntos descansa en dos propiedades matemáticas inquebrantables.

\begin{teoremabox}[Teorema 2.2: Espectro y Completitud de Operadores Hermíticos]
    Sea $\hat{A}$ un operador lineal hermítico ($\hat{A}\adj = \hat{A}$) definido en un espacio de Hilbert $\mathcal{H}$ de dimensión finita $N$ (o un operador autoadjunto compacto en un espacio infinito-dimensional). Entonces:
    \begin{enumerate}
        \item \textbf{Realidad de los Autovalores:} Todos los autovalores de $\hat{A}$ son estrictamente reales:
        \begin{equation}
            \hat{A}\ket{\phi_n} = a_n\ket{\phi_n} \implies a_n \in \mathbb{R}.
        \end{equation}
        \item \textbf{Ortogonalidad de Autoestados:} Los autovectores asociados a autovalores distintos son mutuamente ortogonales:
        \begin{equation}
            a_n \neq a_m \implies \braket{\phi_n}{\phi_m} = 0.
        \end{equation}
        \item \textbf{Completitud y Relación de Clausura:} Los autovectores normalizados $\{\ket{\phi_n}\}_{n=1}^N$ conforman una base ortonormal completa de $\mathcal{H}$, de modo que el operador identidad admite la resolución espectral:
        \begin{equation}
            \sum_{n=1}^N \ket{\phi_n}\bra{\phi_n} = \hat{I}.
        \end{equation}
        \item \textbf{Descomposición Espectral y Diagonalización:} El operador $\hat{A}$ puede expandirse de forma única en términos de sus proyectores ortogonales unidimensionales $\hat{P}_n \equiv \ket{\phi_n}\bra{\phi_n}$:
        \begin{equation}
            \hat{A} = \sum_{n=1}^N a_n \ket{\phi_n}\bra{\phi_n} = \sum_{n=1}^N a_n \hat{P}_n.
            \label{eq:descomposicion_espectral_teorica}
        \end{equation}
        En dicha base propia, la matriz representante de $\hat{A}$ es estrictamente diagonal:
        \begin{equation}
            A_{nm} = \matrixel{\phi_n}{\hat{A}}{\phi_m} = a_n \delta_{nm}.
        \end{equation}
    \end{enumerate}
\end{teoremabox}

\textbf{Demostración Didáctica y Rigurosa:}\\
\textit{1. Demostración de que los autovalores son reales:}
Sea $\ket{\phi_n}$ un autovector normalizado de $\hat{A}$ con autovalor $a_n$:
\begin{equation}
    \hat{A}\ket{\phi_n} = a_n\ket{\phi_n}.
\end{equation}
Multiplicamos por la izquierda por el bra $\bra{\phi_n}$:
\begin{equation}
    \matrixel{\phi_n}{\hat{A}}{\phi_n} = a_n\braket{\phi_n}{\phi_n} = a_n.
\end{equation}
Tomando el complejo conjugado y usando la simetría hermítica del producto interno y la condición $\hat{A}\adj = \hat{A}$:
\begin{equation}
    a_n^* = \left(\matrixel{\phi_n}{\hat{A}}{\phi_n}\right)^* = \matrixel{\phi_n}{\hat{A}\adj}{\phi_n} = \matrixel{\phi_n}{\hat{A}}{\phi_n} = a_n.
\end{equation}
Como $a_n^* = a_n$, el autovalor $a_n$ debe ser forzosamente un número real ($a_n \in \mathbb{R}$).

\textit{2. Demostración de la ortogonalidad para autovalores distintos:}
Consideremos dos autovectores $\ket{\phi_n}$ y $\ket{\phi_m}$ tales que:
\begin{equation}
    \hat{A}\ket{\phi_n} = a_n\ket{\phi_n}, \qquad \hat{A}\ket{\phi_m} = a_m\ket{\phi_m}, \quad \text{con } a_n \neq a_m.
\end{equation}
Evaluemos el elemento de matriz $\matrixel{\phi_m}{\hat{A}}{\phi_n}$ de dos maneras distintas:
\begin{align}
    \text{Por la derecha (acción sobre ket):} \quad &\matrixel{\phi_m}{\hat{A}}{\phi_n} = \bra{\phi_m}(a_n\ket{\phi_n}) = a_n\braket{\phi_m}{\phi_n}. \\
    \text{Por la izquierda (acción sobre bra):} \quad &\matrixel{\phi_m}{\hat{A}}{\phi_n} = \braket{\hat{A}\adj\phi_m}{\phi_n} = \braket{\hat{A}\phi_m}{\phi_n} = (a_m\bra{\phi_m})\ket{\phi_n} = a_m\braket{\phi_m}{\phi_n},
\end{align}
donde hemos utilizado que $a_m^* = a_m$. Restando ambas igualdades:
\begin{equation}
    (a_n - a_m)\braket{\phi_m}{\phi_n} = 0.
\end{equation}
Dado que por hipótesis los autovalores son distintos ($a_n \neq a_m \implies a_n - a_m \neq 0$), el producto interno debe anularse obligatoriamente:
\begin{equation}
    \boxed{\braket{\phi_m}{\phi_n} = 0 \quad (\forall\,n \neq m)}.
\end{equation}

\textit{3. Completitud y diagonalización (inducción sobre la dimensión):}\\
Por el teorema fundamental del álgebra, el polinomio característico de $\hat{A}$ posee al menos una raíz en $\mathbb{C}$, garantizando la existencia de un primer autovector normalizado $\ket{\phi_1}$ con autovalor real $a_1$. Sea $W = \mathrm{span}\{\ket{\phi_1}\}$ y $W^\perp$ su complemento ortogonal en $\mathcal{H}$. Para todo $\ket{\psi} \in W^\perp$, se cumple $\braket{\phi_1}{\hat{A}\psi} = \matrixel{\phi_1}{\hat{A}}{\psi} = \braket{\hat{A}\phi_1}{\psi} = a_1\braket{\phi_1}{\psi} = 0$, lo que demuestra que el subespacio $W^\perp$ de dimensión $N-1$ es invariante bajo $\hat{A}$. Por hipótesis inductiva sobre la dimensión, la restricción de $\hat{A}$ a $W^\perp$ es hermítica y admite una base ortonormal de autovectores. Si un autovalor $a_n$ posee multiplicidad geométrica $g_n > 1$ (subespacio propio degenerado $E_{a_n}$), el procedimiento de ortonormalización de Gram--Schmidt permite elegir siempre una base ortonormal interna para $E_{a_n}$. La unión de dichas bases conforma una base ortonormal completa de todo el espacio $\mathcal{H}$, demostrando la relación de clausura $\sum_{n=1}^N \ket{\phi_n}\bra{\phi_n} = \hat{I}$ y la diagonalizabilidad de $\hat{A}$.

\subsection{Conmutación de Observables y Diagonalización Simultánea (Teorema 2.3)}
\label{subsec:teorema_2_3}

Uno de los resultados más trascendentales de la física cuántica establece la relación íntima entre la \textbf{compatibilidad física de dos mediciones} y la \textbf{conmutación algebraica de sus operadores}.

\begin{teoremabox}[Teorema 2.3: Operadores Compatibles y Base Común]
    Dos operadores hermíticos $\hat{A}$ y $\hat{B}$ admiten una base ortonormal completa de autovectores simultáneos comunes si y solo si \textbf{conmutan entre sí}:
    \begin{equation}
        [\hat{A}, \hat{B}] = \hat{A}\hat{B} - \hat{B}\hat{A} = 0.
    \end{equation}
\end{teoremabox}

\textbf{Demostración de la Condición Necesaria ($\Leftarrow$):}\\
Si suponemos que existe una base ortonormal completa común $\{\ket{n}\}$ tal que $\hat{A}\ket{n} = a_n\ket{n}$ y $\hat{B}\ket{n} = b_n\ket{n}$ para todo $n$, entonces:
\begin{equation}
    [\hat{A}, \hat{B}]\ket{n} = (\hat{A}\hat{B} - \hat{B}\hat{A})\ket{n} = (a_n b_n - b_n a_n)\ket{n} = 0.
    \label{a1-eq:comm_necessary}
\end{equation}
Como los autovalores escalares conmutan ($a_n b_n = b_n a_n$) y $\{\ket{n}\}$ es una base completa que genera todo $\mathcal{H}$, el operador $[\hat{A}, \hat{B}]$ anula cualquier combinación lineal $\ket{\psi} = \sum_n c_n\ket{n}$, concluyéndose idénticamente que $[\hat{A}, \hat{B}] = 0$.

\textbf{Demostración de la Condición Suficiente ($\Rightarrow$):}\\
Para probar que $[\hat{A}, \hat{B}] = 0$ implica la existencia de una base común y prepararnos para resolver problemas con operadores compatibles, debemos distinguir con absoluto rigor didáctico los dos escenarios posibles: el caso sin degeneración y el caso con degeneración.

\subsubsection{Caso 1: Espectro No Degenerado ($g = 1$)}

Supongamos primero que todos los autovalores del operador $\hat{A}$ son simples (no degenerados), es decir, a cada autovalor $a_n$ le corresponde un subespacio propio unidimensional generado por el autovector $\ket{\phi_n}$:
\begin{equation}
    \hat{A}\ket{\phi_n} = a_n\ket{\phi_n}, \quad \text{con } a_n \neq a_m \text{ para } n \neq m.
\end{equation}
Apliquemos el operador $\hat{B}$ sobre el vector $\ket{\phi_n}$ y examinemos qué le hace el operador $\hat{A}$ al vector resultante $\hat{B}\ket{\phi_n}$. Como $[\hat{A}, \hat{B}] = 0$:
\begin{equation}
    \hat{A}(\hat{B}\ket{\phi_n}) = (\hat{A}\hat{B})\ket{\phi_n} = (\hat{B}\hat{A})\ket{\phi_n} = \hat{B}(\hat{A}\ket{\phi_n}) = \hat{B}(a_n\ket{\phi_n}) = a_n(\hat{B}\ket{\phi_n}).
    \label{eq:conmutacion_no_deg}
\end{equation}
¡Observemos con atención la Ec.~\eqref{eq:conmutacion_no_deg}! Nos dice que el vector $\ket{\chi} \equiv \hat{B}\ket{\phi_n}$ es un autovector de $\hat{A}$ con el mismo autovalor $a_n$.\\
Pero como el subespacio propio $E_{a_n}$ asociado al autovalor $a_n$ es unidimensional (no degenerado), \textbf{todo vector en dicho subespacio debe ser forzosamente un múltiplo escalar del vector base $\ket{\phi_n}$}:
\begin{equation}
    \hat{B}\ket{\phi_n} = b_n\ket{\phi_n}, \quad \text{para algún escalar } b_n \in \mathbb{C}.
\end{equation}
Como $\hat{B}$ es hermítico, el escalar $b_n$ es necesariamente real ($b_n \in \mathbb{R}$). Por lo tanto, el vector $\ket{\phi_n}$ es simultáneamente autovector de $\hat{A}$ (con autovalor $a_n$) y de $\hat{B}$ (con autovalor $b_n$):
\begin{equation}
    \boxed{\hat{A}\ket{\phi_n} = a_n\ket{\phi_n} \quad \text{y} \quad \hat{B}\ket{\phi_n} = b_n\ket{\phi_n}}.
\end{equation}
En el caso no degenerado, la base ortonormal común es \textbf{única} (salvo fases globales de cada vector).

\subsubsection{Caso 2: Espectro Degenerado ($g > 1$) y Subespacios Invariantes}

¿Qué ocurre si el operador $\hat{A}$ presenta degeneración espectral? Supongamos que a un determinado autovalor $a_n$ le corresponde un subespacio propio $E_{a_n}$ de dimensión $g_n > 1$ (multiplicidad algebraica/geométrica):
\begin{equation}
    E_{a_n} \equiv \left\{ \ket{\psi} \in \mathcal{H} : \hat{A}\ket{\psi} = a_n\ket{\psi} \right\}, \quad \dim(E_{a_n}) = g_n > 1.
\end{equation}
Sea $\{\ket{\phi_{n, i}}\}_{i=1}^{g_n}$ una base ortonormal cualquiera elegida para dicho subespacio.
\begin{enumerate}
    \item \textbf{Invariancia del subespacio propio bajo $\hat{B}$:}
    Si $\ket{\psi} \in E_{a_n}$, entonces por la conmutación $[\hat{A}, \hat{B}] = 0$:
    \begin{equation}
        \hat{A}(\hat{B}\ket{\psi}) = \hat{B}\hat{A}\ket{\psi} = \hat{B}(a_n\ket{\psi}) = a_n(\hat{B}\ket{\psi}) \implies \hat{B}\ket{\psi} \in E_{a_n}.
    \end{equation}
    Esto demuestra un hecho geométrico fundamental: \textbf{el operador $\hat{B}$ no saca los vectores de $E_{a_n}$; el subespacio $E_{a_n}$ es estrictamente invariante bajo $\hat{B}$}.
    \item \textbf{¿Es $\ket{\phi_{n, i}}$ automáticamente un autovector de $\hat{B}$? ¡NO!}\\
    En general, la acción de $\hat{B}$ sobre un vector base $\ket{\phi_{n, i}}$ es una combinación lineal arbitraria de los vectores que forman la base de $E_{a_n}$:
    \begin{equation}
        \hat{B}\ket{\phi_{n, i}} = \sum_{j=1}^{g_n} B_{ji}^{(n)} \ket{\phi_{n, j}} \neq \text{constante} \times \ket{\phi_{n, i}}.
    \end{equation}
    Por consiguiente, un autovector genérico de $\hat{A}$ dentro del subespacio degenerado \textbf{no es necesariamente un autovector de $\hat{B}$}.
    \item \textbf{Diagonalización interna en el bloque degenerado:}\\
    La restricción de $\hat{B}$ al subespacio $E_{a_n}$ (que denotamos por $\hat{B}|_{E_{a_n}}$) es un operador hermítico que actúa en un espacio de dimensión finita $g_n$. Por el Teorema 2.2, dicho operador restringido \textbf{puede ser diagonalizado de forma independiente dentro de $E_{a_n}$}.\\
    Existen exactamente $g_n$ combinaciones lineales ortonormales $\{\ket{\chi_{n, k}}\}_{k=1}^{g_n} \subset E_{a_n}$ tales que:
    \begin{equation}
        \hat{B}\ket{\chi_{n, k}} = b_{n, k}\ket{\chi_{n, k}}, \quad k \in \{1, 2, \dots, g_n\}.
    \end{equation}
    Y como cada $\ket{\chi_{n, k}}$ pertenece al subespacio $E_{a_n}$, automáticamente satisface:
    \begin{equation}
        \hat{A}\ket{\chi_{n, k}} = a_n \ket{\chi_{n, k}}.
    \end{equation}
    ¡Los vectores $\{\ket{\chi_{n, k}}\}$ son los autovectores comunes buscados!
    \item \textbf{No unicidad de la base:}\\
    Si algunos autovalores $b_{n, k}$ del operador $\hat{B}$ vuelven a coincidir (degeneración en $\hat{B}$), la base común simultánea \textbf{tampoco es única}. Se requerirá añadir un tercer observable compatible $\hat{C}$ ($[\hat{A}, \hat{C}] = [\hat{B}, \hat{C}] = 0$) para terminar de romper la degeneración y etiquetar cada estado de manera no ambigua. ¡Este es exactamente el origen del concepto de \textbf{Conjunto Completo de Observables Compatibles (CCOC)}!
\end{enumerate}

\subsubsection{Ejemplo Matricial Guiado: Diagonalización Simultánea con Degeneración}

Para consolidar la técnica de diagonalización simultánea, desarrollemos paso a paso un problema canónico.

\begin{ejemplobox}[Ejemplo Resuelto: Diagonalización Simultánea de Matrices $3\times 3$ con Degeneración]
    \textbf{Enunciado:} Considere dos operadores lineales representados en la base canónica de $\mathbb{R}^3$ por las matrices simétricas reales:
    \begin{equation}
        A = \begin{pmatrix} 1 & 0 & 0 \\ 0 & 1 & 0 \\ 0 & 0 & -2 \end{pmatrix}, \qquad B = \begin{pmatrix} 0 & 2 & 0 \\ 2 & 0 & 0 \\ 0 & 0 & 3 \end{pmatrix}.
    \end{equation}
    \begin{enumerate}
        \item Verifique explícitamente que $A$ y $B$ conmutan ($[A, B] = 0$).
        \item Halle los autovalores de $A$, identifique el subespacio propio degenerado $E_1$ y explique por qué la base canónica ya no es autovectorial para $B$.
        \item Restrinja $B$ al subespacio degenerado $E_1$, diagonalícelo y construya la matriz unitaria (ortogonal) de cambio de base $P$ que diagonaliza simultáneamente a ambas matrices.
    \end{enumerate}

    \textbf{Solución Paso a Paso:}

    \textit{Paso 1: Verificación de la Conmutación:}
    Efectuamos el producto matricial en ambos órdenes:
    \begin{align}
        AB &= \begin{pmatrix} 1 & 0 & 0 \\ 0 & 1 & 0 \\ 0 & 0 & -2 \end{pmatrix} \begin{pmatrix} 0 & 2 & 0 \\ 2 & 0 & 0 \\ 0 & 0 & 3 \end{pmatrix} = \begin{pmatrix} 0 & 2 & 0 \\ 2 & 0 & 0 \\ 0 & 0 & -6 \end{pmatrix}, \\
        BA &= \begin{pmatrix} 0 & 2 & 0 \\ 2 & 0 & 0 \\ 0 & 0 & 3 \end{pmatrix} \begin{pmatrix} 1 & 0 & 0 \\ 0 & 1 & 0 \\ 0 & 0 & -2 \end{pmatrix} = \begin{pmatrix} 0 & 2 & 0 \\ 2 & 0 & 0 \\ 0 & 0 & -6 \end{pmatrix}.
    \end{align}
    Puesto que $AB = BA$, se verifica idénticamente que $[A, B] = 0$.

    \textit{Paso 2: Espectro de $A$ y Subespacios Propios:}
    Como la matriz $A$ ya es diagonal, sus autovalores se leen directamente de la diagonal principal:
    \begin{equation}
        \lambda_1 = \lambda_2 = 1 \;\; (\text{deg., } g=2), \qquad \lambda_3 = -2 \;\; (\text{no deg., } g=1).
    \end{equation}
    El autovector asociado al autovalor no degenerado $\lambda_3 = -2$ es único (salvo fase):
    \begin{equation}
        \ket{v_3} = \begin{pmatrix} 0 \\ 0 \\ 1 \end{pmatrix}.
    \end{equation}
    El subespacio degenerado asociado al autovalor doble $\lambda = 1$ es el plano generado por los dos primeros vectores de la base canónica:
    \begin{equation}
        E_1 = \mathrm{span}\left\{ \ket{e_1} = \begin{pmatrix} 1 \\ 0 \\ 0 \end{pmatrix}, \, \ket{e_2} = \begin{pmatrix} 0 \\ 1 \\ 0 \end{pmatrix} \right\}.
    \end{equation}
    Notemos con absoluta claridad pedagógica: aunque $\ket{e_1}$ y $\ket{e_2}$ son autovectores de $A$ con autovalor $1$, \textbf{ninguno de ellos es autovector de $B$}, pues:
    \begin{equation}
        B\ket{e_1} = \begin{pmatrix} 0 & 2 & 0 \\ 2 & 0 & 0 \\ 0 & 0 & 3 \end{pmatrix}\begin{pmatrix} 1 \\ 0 \\ 0 \end{pmatrix} = \begin{pmatrix} 0 \\ 2 \\ 0 \end{pmatrix} = 2\ket{e_2} \neq \text{constante} \times \ket{e_1}.
    \end{equation}
    El operador $B$ \enquote{mezcla} los vectores dentro del plano $E_1$.

    \textit{Paso 3: Restricción y Diagonalización en el Bloque Degenerado:}
    La matriz de $B$ restringida al subespacio bidimensional $E_1 = \mathrm{span}\{\ket{e_1}, \ket{e_2}\}$ viene dada por el bloque superior izquierdo de $2\times 2$:
    \begin{equation}
        B|_{E_1} = \begin{pmatrix} 0 & 2 \\ 2 & 0 \end{pmatrix}.
    \end{equation}
    Hallamos sus autovalores resolviendo el polinomio característico secular:
    \begin{equation}
        \det(B|_{E_1} - \mu I_2) = \begin{vmatrix} -\mu & 2 \\ 2 & -\mu \end{vmatrix} = \mu^2 - 4 = 0 \implies \mu = \pm 2.
    \end{equation}
    Calculamos los autovectores correspondientes en el subespacio $E_1$:
    \begin{itemize}
        \item Para $\mu_1 = +2$:
        \begin{equation}
            \begin{pmatrix} -2 & 2 \\ 2 & -2 \end{pmatrix}\begin{pmatrix} c_1 \\ c_2 \end{pmatrix} = \begin{pmatrix} 0 \\ 0 \end{pmatrix} \implies c_1 = c_2 \implies \ket{u_1} = \frac{1}{\sqrt{2}}\begin{pmatrix} 1 \\ 1 \\ 0 \end{pmatrix}.
        \end{equation}
        \item Para $\mu_2 = -2$:
        \begin{equation}
            \begin{pmatrix} 2 & 2 \\ 2 & 2 \end{pmatrix}\begin{pmatrix} c_1 \\ c_2 \end{pmatrix} = \begin{pmatrix} 0 \\ 0 \end{pmatrix} \implies c_1 = -c_2 \implies \ket{u_2} = \frac{1}{\sqrt{2}}\begin{pmatrix} 1 \\ -1 \\ 0 \end{pmatrix}.
        \end{equation}
    \end{itemize}
    Notemos que $\ket{u_1}$ y $\ket{u_2}$ siguen perteneciendo a $E_1$, por lo que $A\ket{u_1} = 1\ket{u_1}$ y $A\ket{u_2} = 1\ket{u_2}$.\\
    Para el tercer autovector, evaluamos la acción de $B$ sobre $\ket{v_3}$:
    \begin{equation}
        B\ket{v_3} = \begin{pmatrix} 0 & 2 & 0 \\ 2 & 0 & 0 \\ 0 & 0 & 3 \end{pmatrix}\begin{pmatrix} 0 \\ 0 \\ 1 \end{pmatrix} = \begin{pmatrix} 0 \\ 0 \\ 3 \end{pmatrix} = 3\ket{v_3}.
    \end{equation}
    Por tanto, $\ket{u_3} \equiv \ket{v_3} = \begin{pmatrix} 0 \\ 0 \\ 1 \end{pmatrix}$ ya era autovector de $B$ con autovalor $\mu_3 = 3$.

    \textit{Paso 4: Construcción de la Matriz de Paso $P$ y Verificación Final:}
    Colocamos los autovectores comunes $\{\ket{u_1}, \ket{u_2}, \ket{u_3}\}$ como columnas de la matriz ortogonal de transformación $P$:
    \begin{equation}
        P = \begin{pmatrix}
            \frac{1}{\sqrt{2}} & \frac{1}{\sqrt{2}} & 0 \\
            \frac{1}{\sqrt{2}} & -\frac{1}{\sqrt{2}} & 0 \\
            0 & 0 & 1
        \end{pmatrix}.
    \end{equation}
    Puesto que las columnas son ortonormales, $P$ es ortogonal ($P^{-1} = P^T$). Comprobemos que diagonaliza simultáneamente a ambas matrices:
    \begin{align}
        P^{-1} A P &= P^T A P = \begin{pmatrix}
            \frac{1}{\sqrt{2}} & \frac{1}{\sqrt{2}} & 0 \\
            \frac{1}{\sqrt{2}} & -\frac{1}{\sqrt{2}} & 0 \\
            0 & 0 & 1
        \end{pmatrix} \begin{pmatrix} 1 & 0 & 0 \\ 0 & 1 & 0 \\ 0 & 0 & -2 \end{pmatrix} \begin{pmatrix}
            \frac{1}{\sqrt{2}} & \frac{1}{\sqrt{2}} & 0 \\
            \frac{1}{\sqrt{2}} & -\frac{1}{\sqrt{2}} & 0 \\
            0 & 0 & 1
        \end{pmatrix} = \begin{pmatrix} 1 & 0 & 0 \\ 0 & 1 & 0 \\ 0 & 0 & -2 \end{pmatrix}, \\
        P^{-1} B P &= P^T B P = \begin{pmatrix}
            \frac{1}{\sqrt{2}} & \frac{1}{\sqrt{2}} & 0 \\
            \frac{1}{\sqrt{2}} & -\frac{1}{\sqrt{2}} & 0 \\
            0 & 0 & 1
        \end{pmatrix} \begin{pmatrix} 0 & 2 & 0 \\ 2 & 0 & 0 \\ 0 & 0 & 3 \end{pmatrix} \begin{pmatrix}
            \frac{1}{\sqrt{2}} & \frac{1}{\sqrt{2}} & 0 \\
            \frac{1}{\sqrt{2}} & -\frac{1}{\sqrt{2}} & 0 \\
            0 & 0 & 1
        \end{pmatrix} = \begin{pmatrix} 2 & 0 & 0 \\ 0 & -2 & 0 \\ 0 & 0 & 3 \end{pmatrix}.
    \end{align}
    ¡Ambas matrices quedan rigurosamente diagonalizadas al mismo tiempo en la base $\{\ket{u_1}, \ket{u_2}, \ket{u_3}\}$!
\end{ejemplobox}

\begin{figure}[H]
\centering
\begin{tikzpicture}[
    x={(2.1cm,0cm)}, y={(1.0cm,0.8cm)}, z={(0cm,2.0cm)},
    >=Stealth, font=\small]

  % plano degenerado E_1
  \fill[azulUNI!10] (-1.1,-1.1,0)--(1.1,-1.1,0)--(1.1,1.1,0)--(-1.1,1.1,0)--cycle;
  \draw[azulUNI!45] (-1.1,-1.1,0)--(1.1,-1.1,0)--(1.1,1.1,0)--(-1.1,1.1,0)--cycle;
  \node[azulUNI,anchor=south east,font=\small] at (-0.85,1.1,0)
        {$E_1$ \; ($A=1$, degenerado)};

  % base canónica (no es propia de B)
  \draw[->,thick,grisUNI] (0,0,0)--(1,0,0) node[below right]{$|e_1\rangle$};
  \draw[->,thick,grisUNI] (0,0,0)--(0,1,0) node[above left]{$|e_2\rangle$};

  % base propia común
  \draw[->,very thick,rojoUNI]  (0,0,0)--(0.85,0.85,0)
        node[above right]{$|u_1\rangle$};
  \draw[->,very thick,azulUNI]  (0,0,0)--(0.85,-0.85,0)
        node[below]{$|u_2\rangle$};
  \draw[->,very thick,verdeUNI] (0,0,0)--(0,0,1.4)
        node[above]{$|v_3\rangle$};

  % B mezcla e_1 y e_2
  \draw[<->,dashed,grisUNI] (0.55,0,0) to[bend left=35] (0,0.55,0);

  % leyenda
  \node[align=left,anchor=north west,font=\small] at (canvas cs:x=-3.6cm,y=-1.7cm) {%
    \textcolor{rojoUNI}{$|u_1\rangle=\tfrac{1}{\sqrt2}(|e_1\rangle+|e_2\rangle)$: $A=1,\;B=+2$}\\[2pt]
    \textcolor{azulUNI}{$|u_2\rangle=\tfrac{1}{\sqrt2}(|e_1\rangle-|e_2\rangle)$: $A=1,\;B=-2$}\\[2pt]
    \textcolor{verdeUNI}{$|v_3\rangle=|e_3\rangle$: $A=-2,\;B=3$}\\[6pt]
    \textcolor{grisUNI}{$B|e_1\rangle=2|e_2\rangle$\; (\,$B$ «mezcla» dentro de $E_1$)}};
\end{tikzpicture}
\caption{Diagonalización simultánea con degeneración (Ejemplo de la Sección~2.8.3). Dentro del
plano $E_1$, $A$ es un múltiplo de la identidad, pero $B$ gira los vectores:
la base canónica $\{|e_1\rangle,|e_2\rangle\}$ no es propia de $B$. La base
común se obtiene diagonalizando $B$ restringido a $E_1$, lo que equivale a
girar $45^\circ$ dentro del plano; $|v_3\rangle$ ya era propio de ambos.}
\label{fig:diag-simultanea}
\end{figure}


\subsubsection{Ejemplos Físicos de Conjuntos Simultáneamente Diagonalizables}

Este procedimiento algebraico se manifiesta en sistemas físicos fundamentales de la física cuántica:
\begin{enumerate}
    \item \textbf{Partícula Libre en 3D:} Las tres componentes cartesianas del momento lineal conmutan entre sí:
    \begin{equation}
        [\hat{p}_x, \hat{p}_y] = [\hat{p}_y, \hat{p}_z] = [\hat{p}_z, \hat{p}_x] = 0, \quad \text{y además } [\hat{H}, \hat{p}_i] = 0 \quad \text{con } \hat{H} = \frac{\hat{p}^2}{2m}.
    \end{equation}
    Admiten como base común continua de autoestados a las \textbf{ondas planas tridimensionales}:
    \begin{equation}
        \psi_{\vec{p}}(\vec{r}) = \braket{\vec{r}}{\vec{p}} = \frac{1}{(2\pi\hbar)^{3/2}}\eu^{\frac{\iu}{\hbar}\vec{p}\cdot\vec{r}}, \quad \text{con } \hat{\vec{p}}\ket{\vec{p}} = \vec{p}\ket{\vec{p}} \quad \text{y} \quad \hat{H}\ket{\vec{p}} = \frac{p^2}{2m}\ket{\vec{p}}.
    \end{equation}
    \item \textbf{Oscilador Armónico 3D Isótropo Separable:} El Hamiltoniano total es la suma de tres osciladores 1D independientes:
    \begin{equation}
        \hat{H} = \hat{H}_x + \hat{H}_y + \hat{H}_z = \hbar\omega\left(\hat{a}_x\adj\hat{a}_x + \hat{a}_y\adj\hat{a}_y + \hat{a}_z\adj\hat{a}_z + \frac{3}{2}\right).
    \end{equation}
    Dado que $[\hat{H}_x, \hat{H}_y] = [\hat{H}_y, \hat{H}_z] = [\hat{H}_z, \hat{H}_x] = 0$, la base ortonormal común viene dada por los productos directos de autoestados de Fock unidimensionales:
    \begin{equation}
        \ket{n_x, n_y, n_z} = \ket{n_x} \otimes \ket{n_y} \otimes \ket{n_z}, \quad \text{con energía } E = \hbar\omega\left(n_x + n_y + n_z + \frac{3}{2}\right).
    \end{equation}
\end{enumerate}

\subsection{Taller de Conmutadores Fundamentales, Par Canónico y Operadores Escalera}
\label{subsec:taller_conmutadores}

Antes de abordar las demostraciones avanzadas de las secciones siguientes, resulta indispensable dominar la \enquote{carpintería algebraica} de los conmutadores, comprender por qué la posición y el momento no conmutan, y conocer el álgebra de los operadores de creación y aniquilación.

\subsubsection{Las Cuatro Reglas Operativas del Álgebra de Conmutadores}

Dados tres operadores lineales cualesquiera $\hat{A}, \hat{B}, \hat{C}$ y escalares $\alpha, \beta \in \mathbb{C}$, el conmutador $[\hat{A}, \hat{B}] \equiv \hat{A}\hat{B} - \hat{B}\hat{A}$ satisface:
\begin{enumerate}
    \item \textbf{Antisimetría:} Intercambiar el orden de los operadores invierte el signo:
    \begin{equation}
        [\hat{A}, \hat{B}] = -[\hat{B}, \hat{A}], \quad \text{en particular } [\hat{A}, \hat{A}] = 0 \text{ y } [\hat{A}, c\hat{I}] = 0.
    \end{equation}
    \item \textbf{Bilinealidad:} Es lineal respecto a la suma y multiplicación por escalares en ambos argumentos:
    \begin{equation}
        [\alpha\hat{A} + \beta\hat{B}, \, \hat{C}] = \alpha[\hat{A}, \hat{C}] + \beta[\hat{B}, \hat{C}], \qquad [\hat{A}, \, \alpha\hat{B} + \beta\hat{C}] = \alpha[\hat{A}, \hat{B}] + \beta[\hat{A}, \hat{C}].
    \end{equation}
    \item \textbf{Regla de Leibniz (Derivación de un Producto):} Permite descomponer el conmutador de un producto en conmutadores simples:
    \begin{equation}
        \boxed{[\hat{A}\hat{B}, \hat{C}] = \hat{A}[\hat{B}, \hat{C}] + [\hat{A}, \hat{C}]\hat{B}}, \qquad
        \boxed{[\hat{A}, \hat{B}\hat{C}] = [\hat{A}, \hat{B}]\hat{C} + \hat{B}[\hat{A}, \hat{C}]}.
        \label{eq:leibniz_comm_basic}
    \end{equation}
    \textit{Demostración directa:} Expandiendo el miembro derecho de la primera identidad y observando la cancelación telescópica del término intermedio:
    \begin{equation}
        \hat{A}[\hat{B}, \hat{C}] + [\hat{A}, \hat{C}]\hat{B} = \hat{A}\hat{B}\hat{C} - \bcancel{\hat{A}\hat{C}\hat{B}} + \bcancel{\hat{A}\hat{C}\hat{B}} - \hat{C}\hat{A}\hat{B} = \hat{A}\hat{B}\hat{C} - \hat{C}\hat{A}\hat{B} = [\hat{A}\hat{B}, \hat{C}].
    \end{equation}
    \item \textbf{Identidad de Jacobi:} La suma cíclica de conmutadores anidados es idénticamente nula (propiedad definitoria de toda álgebra de Lie):
    \begin{equation}
        [\hat{A}, [\hat{B}, \hat{C}]] + [\hat{B}, [\hat{C}, \hat{A}]] + [\hat{C}, [\hat{A}, \hat{B}]] = 0.
    \end{equation}
\end{enumerate}

\subsubsection{¿De dónde sale \texorpdfstring{$[\hat{x}, \hat{p}_x] = \iu\hbar\hat{I}$}{[x, px] = ihbar I}? Demostración sobre una Función de Prueba}

En la representación de coordenadas espaciales, el operador posición actúa multiplicando por $x$ ($\hat{x}\psi(x) = x\psi(x)$) y el operador momento lineal actúa derivando ($\hat{p}_x\psi(x) = -\iu\hbar\frac{\dif\psi(x)}{\dif x}$).

\begin{alertabox}[Regla de Oro para Evaluar Conmutadores Diferenciales]
    Nunca evalúe un operador diferencial \enquote{en el vacío}. Para hallar $[\hat{x}, \hat{p}_x]$, hágalo actuar sobre una función de onda diferenciable arbitraria (\textbf{función de prueba}) $\psi(x)$ y aplique la regla de la derivada de un producto:
    \begin{align}
        [\hat{x}, \hat{p}_x]\psi(x) &= (\hat{x}\hat{p}_x - \hat{p}_x\hat{x})\psi(x) = x\left(-\iu\hbar\frac{\dif\psi}{\dif x}\right) - \left(-\iu\hbar\frac{\dif}{\dif x}\right)\big(x\psi(x)\big) \nonumber\\
        &= -\iu\hbar\,x\frac{\dif\psi}{\dif x} + \iu\hbar\left( \frac{\dif x}{\dif x}\psi(x) + x\frac{\dif\psi}{\dif x} \right) \nonumber\\
        &= -\bcancel{\iu\hbar\,x\frac{\dif\psi}{\dif x}} + \iu\hbar\,\psi(x) + \bcancel{\iu\hbar\,x\frac{\dif\psi}{\dif x}} = \iu\hbar\,\psi(x).
    \end{align}
    Como esta igualdad se verifica para cualquier función $\psi(x) \in \mathcal{H}$, suprimiendo la función de prueba a ambos lados se obtiene la célebre \textbf{Relación de Conmutación Canónica de Born--Jordan--Heisenberg}:
    \begin{equation}
        \boxed{[\hat{x}, \hat{p}_x] = \iu\hbar\hat{I}}, \quad \text{y en 3D:} \quad \boxed{[\hat{x}_i, \hat{p}_j] = \iu\hbar\,\delta_{ij}\hat{I}}.
    \end{equation}
\end{alertabox}

\subsubsection{Operadores Escalera del Oscilador Armónico (\texorpdfstring{$\hat{a}, \hat{a}\adj$}{a, a+}) y Base de Fock}

En múltiples problemas cuánticos (como el Oscilador 2D del Problema Resuelto 5), el Hamiltoniano armónico unidimensional $\hat{H} = \frac{\hat{p}^2}{2m} + \frac{1}{2}m\omega^2\hat{x}^2$ se factoriza algebraicamente definiendo los operadores adimensionales de \textbf{aniquilación} ($\hat{a}$) y \textbf{creación} ($\hat{a}\adj$):
\begin{equation}
    \hat{a} \equiv \sqrt{\frac{m\omega}{2\hbar}}\left(\hat{x} + \frac{\iu}{m\omega}\hat{p}\right), \qquad
    \hat{a}\adj \equiv \sqrt{\frac{m\omega}{2\hbar}}\left(\hat{x} - \frac{\iu}{m\omega}\hat{p}\right).
\end{equation}
Utilizando $[\hat{x}, \hat{p}] = \iu\hbar\hat{I}$, su álgebra fundamental se resume en cuatro relaciones esenciales:
\begin{enumerate}
    \item \textbf{Conmutador bosónico fundamental:} $[\hat{a}, \hat{a}\adj] = \hat{I}$ (y para ejes distintos en 2D/3D, $[\hat{a}_i, \hat{a}_j\adj] = \delta_{ij}\hat{I}$, $[\hat{a}_i, \hat{a}_j] = 0$).
    \item \textbf{Operador Número y Hamiltoniano:} Definiendo el operador hermítico número de cuantos $\hat{N} \equiv \hat{a}\adj\hat{a}$, el Hamiltoniano se reescribe como $\hat{H} = \hbar\omega\left(\hat{N} + \frac{1}{2}\hat{I}\right)$, con conmutadores $[\hat{N}, \hat{a}] = -\hat{a}$ y $[\hat{N}, \hat{a}\adj] = +\hat{a}\adj$.
    \item \textbf{Acción sobre los autoestados de Fock $\{\ket{n}\}$ ($\hat{N}\ket{n} = n\ket{n}$, $n=0,1,2,\dots$):}
    \begin{equation}
        \boxed{\hat{a}\ket{n} = \sqrt{n}\,\ket{n-1} \quad (\text{con } \hat{a}\ket{0} = 0), \qquad \hat{a}\adj\ket{n} = \sqrt{n+1}\,\ket{n+1}}.
    \end{equation}
\end{enumerate}
